\documentclass[10pt,a4paper]{amsart}
\usepackage[margin=19mm]{geometry}
\usepackage{amsmath,amssymb,mathtools}
\usepackage[scr=rsfs,cal=euler]{mathalfa}
\usepackage{xfrac}
\usepackage[unicode,hidelinks,bookmarks=false]{hyperref}
\newcommand{\ie}{i.e.\ }
\newcommand{\cf}{cf.\ }
\newcommand{\resp}{respectively\ }
\newcommand{\Subsection}{\subsection}
\providecommand{\nobreakdash}{\nobreak\hbox{-}}
\setlength{\emergencystretch}{2em}
\multlinegap0pt
\usepackage{amssymb}
\usepackage{mathtools}
\usepackage{bm}
\usepackage{xcolor}
\newtheorem{theorem}{Theorem}
\title{Variational Dissipative Mechanics on Lie Algebroids}
\author{Alexandre Anahory Simoes, Leonardo Colombo}
\date{Source: arXiv:2512.17424v1 (CC BY 4.0)}
\begin{document}
\maketitle

\noindent\emph{Source and licence.} Variational Dissipative Mechanics on Lie Algebroids. Version arXiv:2512.17424v1 (CC BY 4.0), distributed by arXiv under the Creative Commons Attribution 4.0 International licence, http://creativecommons.org/licenses/by/4.0/, as recorded in the arXiv licence metadata for this exact version and shown on the abstract page https://arxiv.org/abs/2512.17424v1. Full licence notice: \texttt{licenses/arxiv-2512.17424-CC-BY-4.0.txt}.\par\smallskip

\noindent\textit{Editorial scope.} Counted unit: the article's theorem thm:ELH\_algebroid (source0.tex lines 352--372). The article's complete original proof of it is delivered with it (source0.tex lines 374--513, one \texttt{proof} environment of 140 lines). No mathematics is written, repaired or re-derived: the statement and the proof are the article's own lines, in the article's own notation, and no retained line is substituted or re-flowed.  Retained uncounted as the source-local context the proof needs: lines 342--344.  Nothing was omitted from any retained span, so the package carries no omission record.  No retained line contains a citation, so the wrapper carries no bibliography and no bibliography entry.  No retained line contains a figure environment, an image-inclusion command or drawing code, so no image is delivered and the package declares no asset dependency.  Editorial formatting only, in addition: an AMS \texttt{amsart} preamble that restates the article's own declarations and macros used by the retained lines, namely the environments \texttt{theorem} on separate counters, together with the standard packages those lines need.  The retained environments print in this document as Theorem 1 in order of appearance, whereas the article prints the same items with the numbering of its own full document.  The complete native source of the article as distributed in the arXiv e-print for this exact version is bundled unchanged as \texttt{sources/191313/source0.tex}.
\begin{equation}\label{eq:stationary}
\left.\frac{d}{ds}\right|_{s=0} z_s(T) = 0.
\end{equation}

\begin{theorem}[Euler--Lagrange--Herglotz equations on a Lie algebroid]
\label{thm:ELH_algebroid}
Let $(E\to M,[\cdot,\cdot],\rho)$ be a Lie algebroid whose anchor components are given by
$\rho^i_{\ \alpha}(x)$ and strcuture functions are $C^\gamma_{\ \alpha\beta}(x)$ with respect to local coordinates denoted by $(x^{i},y^{\alpha},z)$.
Let $L\colon E\times\mathbb{R}\to\mathbb{R}$ be a Herglotz-type Lagrangian and consider a curve $(x^{i}(t),y^{\alpha}(t),z(t))$ satisfying
\begin{equation}\label{eq:ELH_constraints}
\dot x^i(t) = \rho^i_{\ \alpha}(x^{i}(t))\,y^\alpha(t), 
\qquad
\dot z(t) = L(x^{i}(t),y^{\alpha}(t),z(t)).
\end{equation}
Then $(x^{i}(t),y^{\alpha}(t),z(t))$ is a Herglotz extremal, i.e. satisfies
\eqref{eq:stationary} for all admissible variations with $x^{i}(0),x^{i}(T),z(0)$ fixed, 
if and only if it satisfies the Euler--Lagrange--Herglotz equations:
\begin{equation}\label{eq:ELH_equations}
\frac{d}{dt}\left(\frac{\partial L}{\partial y^\alpha}\right)
+ C^\gamma_{\ \alpha\beta}(x)\,y^\beta \frac{\partial L}{\partial y^\gamma}
- \rho^i_{\ \alpha}(x)\,\frac{\partial L}{\partial x^i}
- \frac{\partial L}{\partial z}\,\frac{\partial L}{\partial y^\alpha}
= 0.
\end{equation}
\end{theorem}

\begin{proof}
Let $(x^{i}_s(t),y^{\alpha}_s(t),z_s(t))$ be an admissible variation, i.e.
for each $s$, $\dot x_s^i(t) = \rho^i_{\ \alpha}(x_s(t))\,y_s^\alpha(t)$ and $
\dot z_s(t) = L(x^{i}_s(t),y^{\alpha}_s(t),z_s(t))$, with $x^{i}_s(0)=x^{i}(0)$, $x^{i}_s(T)=x^{i}(T)$, $z_s(0)=z(0)$.
Define $\delta x^i = \partial_s x_s^i|_{s=0}$, 
$\delta y^\alpha=\partial_s y_s^\alpha|_{s=0}$, 
$\delta z = \partial_s z_s|_{s=0}$. Differentiating the Herglotz equation with respect to $s$:
\[
\frac{d}{dt}(\delta z)
= \frac{\partial L}{\partial x^i}\,\delta x^i
+ \frac{\partial L}{\partial y^\alpha}\,\delta y^\alpha
+ \frac{\partial L}{\partial z}\,\delta z.
\]
This is a linear ODE in $\delta z$. Rearranging:
\[
\frac{d}{dt}(\delta z) - \frac{\partial L}{\partial z}\,\delta z
= \frac{\partial L}{\partial x^i}\,\delta x^i
+ \frac{\partial L}{\partial y^\alpha}\,\delta y^\alpha.
\]
Let
\[
\lambda(t) := \exp\!\left(-\int_0^t \frac{\partial L}{\partial z}(\omega)\,d\omega\right).
\]
Then
\[
\frac{d}{dt} \big(\lambda \delta z\big)
= \lambda \left[\frac{d}{dt}(\delta z) - \frac{\partial L}{\partial z}\,\delta z\right]
= \lambda \left(
\frac{\partial L}{\partial x^i}\,\delta x^i
+ \frac{\partial L}{\partial y^\alpha}\,\delta y^\alpha
\right).
\]
Integrating both sides of the previous equation from $0$ to $T$ and using $\delta z(0)=0$ (since $z(0)$ is fixed),
we get
\[
\lambda(T)\,\delta z(T)
= \int_0^T \lambda(t)
\left(
\frac{\partial L}{\partial x^i}\,\delta x^i
+ \frac{\partial L}{\partial y^\alpha}\,\delta y^\alpha
\right) \,dt.
\]
Since $\lambda(T)\neq 0$, the stationarity condition 
$\delta z(T)=0$ is equivalent to
\begin{equation}\label{eq:stationary_integral}
\int_0^T \lambda(t)
\left(
\frac{\partial L}{\partial x^i}\,\delta x^i
+ \frac{\partial L}{\partial y^\alpha}\,\delta y^\alpha
\right) \,dt = 0.
\end{equation}

Now, we express $(\delta x^{i},\delta y^{\alpha})$ in terms of a time-dependent
section $\xi(t) = \xi^\alpha(t)\,e_\alpha|_{x(t)}$ along $x(t)$, via:
\[
\delta x^i(t) = \rho^i_{\ \alpha}(x^{i}(t))\,\xi^\alpha(t),
\qquad
\delta y^\gamma (t)= \dot \xi^\gamma (t) + C^\gamma_{\ \alpha\beta}(x^{i}(t))\,y^\alpha (t) \,\xi^\beta (t).
\]
Substituting into \eqref{eq:stationary_integral}, we obtain
\[
\int_0^T \lambda(t)
\left[
\frac{\partial L}{\partial x^i}\,\rho^i_{\ \alpha}\,\xi^\alpha
+ \frac{\partial L}{\partial y^\gamma}
\left(
\dot \xi^\gamma + C^\gamma_{\ \alpha\beta}y^\alpha\,\xi^\beta
\right)
\right] \,dt = 0.
\]
where we drop the functions inputs to simplify the reading. Rewriting:
\[
\int_0^T \left[
\lambda \frac{\partial L}{\partial x^i}\,\rho^i_{\ \alpha}\,\xi^\alpha
+ \lambda \frac{\partial L}{\partial y^\gamma} C^\gamma_{\ \alpha\beta}y^\alpha\,\xi^\beta
+ \lambda \frac{\partial L}{\partial y^\gamma} \dot \xi^\gamma
\right] dt = 0.
\]

We integrate by parts the term with $\dot \xi^\gamma$:
\[
\int_0^T \lambda \frac{\partial L}{\partial y^\gamma} \dot \xi^\gamma\, dt
= \left[ \lambda \frac{\partial L}{\partial y^\gamma} \xi^\gamma \right]_0^T
- \int_0^T \frac{d}{dt}\Big(\lambda \frac{\partial L}{\partial y^\gamma}\Big)
\,\xi^\gamma \,dt.
\]
The boundary term vanishes because $\xi^{\gamma}(0)=\xi^{\gamma}(T)=0$ (fixed endpoints).
Hence the stationarity condition becomes
\[
\int_0^T \left\{
\lambda \frac{\partial L}{\partial x^i}\,\rho^i_{\ \alpha}
- \lambda \frac{\partial L}{\partial y^\gamma} C^\gamma_{\ \alpha\beta}y^\beta
- \frac{d}{dt}\Big(\lambda \frac{\partial L}{\partial y^\alpha}\Big)
\right\} \xi^\alpha \, dt = 0.
\]
where the middle term changes sign due to the skew-symmetry of the structure functions, i.e., $C_{\alpha \beta}^{\gamma}=-C_{\beta \alpha}^{\gamma}$. Since $\xi^\alpha(t)$ are arbitrary functions vanishing at $t=0,T$,
we deduce the pointwise condition
\[
- \frac{d}{dt}\Big(\lambda \frac{\partial L}{\partial y^\alpha}\Big)
+ \lambda \frac{\partial L}{\partial x^i}\,\rho^i_{\ \alpha}
- \lambda \frac{\partial L}{\partial y^\gamma} C^\gamma_{\ \alpha\beta}y^\beta
= 0.
\]
Expanding the derivative:
\[
\frac{d}{dt}\Big(\lambda \frac{\partial L}{\partial y^\alpha}\Big)
= \dot \lambda \frac{\partial L}{\partial y^\alpha}
+ \lambda \frac{d}{dt}\left( \frac{\partial L}{\partial y^\alpha}\right).
\]
and using $\dot\lambda = -(\partial L/\partial z)\lambda$ by definition of $\lambda$,
we obtain
\[
\dot \lambda \frac{\partial L}{\partial y^\alpha}
= -\lambda \frac{\partial L}{\partial z}\frac{\partial L}{\partial y^\alpha}.
\]
Hence,
\[
- \left[
-\lambda \frac{\partial L}{\partial z}\frac{\partial L}{\partial y^\alpha}
+ \lambda \frac{d}{dt}\left( \frac{\partial L}{\partial y^\alpha}\right)
\right]
+ \lambda \frac{\partial L}{\partial x^i}\,\rho^i_{\ \alpha}
- \lambda \frac{\partial L}{\partial y^\gamma} C^\gamma_{\ \alpha\beta}y^\beta
= 0.
\]
Simplifying and dividing by $\lambda$, which is a nowhere vanishing function, we conclude
\[
\frac{d}{dt}\left(\frac{\partial L}{\partial y^\alpha}\right)
+ C^\gamma_{\ \alpha\beta}(x)\,y^\beta \frac{\partial L}{\partial y^\gamma}
- \rho^i_{\ \alpha}(x)\,\frac{\partial L}{\partial x^i}
- \frac{\partial L}{\partial z}\,\frac{\partial L}{\partial y^\alpha}
= 0.
\]
These are the desired Euler--Lagrange--Herglotz equations \eqref{eq:ELH_equations}.

The converse (that any solution of \eqref{eq:ELH_constraints}--\eqref{eq:ELH_equations}
is stationary) follows by reversing the above computations, showing that then
the integral \eqref{eq:stationary_integral} vanishes for all admissible
variations.
\end{proof}

\end{document}
